% Options for packages loaded elsewhere
% Options for packages loaded elsewhere
\PassOptionsToPackage{unicode}{hyperref}
\PassOptionsToPackage{hyphens}{url}
\PassOptionsToPackage{dvipsnames,svgnames,x11names}{xcolor}
%
\documentclass[
  english,
  letterpaper,
  DIV=11,
  numbers=noendperiod]{scrreprt}
\usepackage{xcolor}
\usepackage{amsmath,amssymb}
\setcounter{secnumdepth}{5}
\usepackage{iftex}
\ifPDFTeX
  \usepackage[T1]{fontenc}
  \usepackage[utf8]{inputenc}
  \usepackage{textcomp} % provide euro and other symbols
\else % if luatex or xetex
  \usepackage{unicode-math} % this also loads fontspec
  \defaultfontfeatures{Scale=MatchLowercase}
  \defaultfontfeatures[\rmfamily]{Ligatures=TeX,Scale=1}
\fi
\usepackage{lmodern}
\ifPDFTeX\else
  % xetex/luatex font selection
\fi
% Use upquote if available, for straight quotes in verbatim environments
\IfFileExists{upquote.sty}{\usepackage{upquote}}{}
\IfFileExists{microtype.sty}{% use microtype if available
  \usepackage[]{microtype}
  \UseMicrotypeSet[protrusion]{basicmath} % disable protrusion for tt fonts
}{}
\makeatletter
\@ifundefined{KOMAClassName}{% if non-KOMA class
  \IfFileExists{parskip.sty}{%
    \usepackage{parskip}
  }{% else
    \setlength{\parindent}{0pt}
    \setlength{\parskip}{6pt plus 2pt minus 1pt}}
}{% if KOMA class
  \KOMAoptions{parskip=half}}
\makeatother
% Make \paragraph and \subparagraph free-standing
\makeatletter
\ifx\paragraph\undefined\else
  \let\oldparagraph\paragraph
  \renewcommand{\paragraph}{
    \@ifstar
      \xxxParagraphStar
      \xxxParagraphNoStar
  }
  \newcommand{\xxxParagraphStar}[1]{\oldparagraph*{#1}\mbox{}}
  \newcommand{\xxxParagraphNoStar}[1]{\oldparagraph{#1}\mbox{}}
\fi
\ifx\subparagraph\undefined\else
  \let\oldsubparagraph\subparagraph
  \renewcommand{\subparagraph}{
    \@ifstar
      \xxxSubParagraphStar
      \xxxSubParagraphNoStar
  }
  \newcommand{\xxxSubParagraphStar}[1]{\oldsubparagraph*{#1}\mbox{}}
  \newcommand{\xxxSubParagraphNoStar}[1]{\oldsubparagraph{#1}\mbox{}}
\fi
\makeatother

\usepackage{color}
\usepackage{fancyvrb}
\newcommand{\VerbBar}{|}
\newcommand{\VERB}{\Verb[commandchars=\\\{\}]}
\DefineVerbatimEnvironment{Highlighting}{Verbatim}{commandchars=\\\{\}}
% Add ',fontsize=\small' for more characters per line
\usepackage{framed}
\definecolor{shadecolor}{RGB}{241,243,245}
\newenvironment{Shaded}{\begin{snugshade}}{\end{snugshade}}
\newcommand{\AlertTok}[1]{\textcolor[rgb]{0.68,0.00,0.00}{#1}}
\newcommand{\AnnotationTok}[1]{\textcolor[rgb]{0.37,0.37,0.37}{#1}}
\newcommand{\AttributeTok}[1]{\textcolor[rgb]{0.40,0.45,0.13}{#1}}
\newcommand{\BaseNTok}[1]{\textcolor[rgb]{0.68,0.00,0.00}{#1}}
\newcommand{\BuiltInTok}[1]{\textcolor[rgb]{0.00,0.23,0.31}{#1}}
\newcommand{\CharTok}[1]{\textcolor[rgb]{0.13,0.47,0.30}{#1}}
\newcommand{\CommentTok}[1]{\textcolor[rgb]{0.37,0.37,0.37}{#1}}
\newcommand{\CommentVarTok}[1]{\textcolor[rgb]{0.37,0.37,0.37}{\textit{#1}}}
\newcommand{\ConstantTok}[1]{\textcolor[rgb]{0.56,0.35,0.01}{#1}}
\newcommand{\ControlFlowTok}[1]{\textcolor[rgb]{0.00,0.23,0.31}{\textbf{#1}}}
\newcommand{\DataTypeTok}[1]{\textcolor[rgb]{0.68,0.00,0.00}{#1}}
\newcommand{\DecValTok}[1]{\textcolor[rgb]{0.68,0.00,0.00}{#1}}
\newcommand{\DocumentationTok}[1]{\textcolor[rgb]{0.37,0.37,0.37}{\textit{#1}}}
\newcommand{\ErrorTok}[1]{\textcolor[rgb]{0.68,0.00,0.00}{#1}}
\newcommand{\ExtensionTok}[1]{\textcolor[rgb]{0.00,0.23,0.31}{#1}}
\newcommand{\FloatTok}[1]{\textcolor[rgb]{0.68,0.00,0.00}{#1}}
\newcommand{\FunctionTok}[1]{\textcolor[rgb]{0.28,0.35,0.67}{#1}}
\newcommand{\ImportTok}[1]{\textcolor[rgb]{0.00,0.46,0.62}{#1}}
\newcommand{\InformationTok}[1]{\textcolor[rgb]{0.37,0.37,0.37}{#1}}
\newcommand{\KeywordTok}[1]{\textcolor[rgb]{0.00,0.23,0.31}{\textbf{#1}}}
\newcommand{\NormalTok}[1]{\textcolor[rgb]{0.00,0.23,0.31}{#1}}
\newcommand{\OperatorTok}[1]{\textcolor[rgb]{0.37,0.37,0.37}{#1}}
\newcommand{\OtherTok}[1]{\textcolor[rgb]{0.00,0.23,0.31}{#1}}
\newcommand{\PreprocessorTok}[1]{\textcolor[rgb]{0.68,0.00,0.00}{#1}}
\newcommand{\RegionMarkerTok}[1]{\textcolor[rgb]{0.00,0.23,0.31}{#1}}
\newcommand{\SpecialCharTok}[1]{\textcolor[rgb]{0.37,0.37,0.37}{#1}}
\newcommand{\SpecialStringTok}[1]{\textcolor[rgb]{0.13,0.47,0.30}{#1}}
\newcommand{\StringTok}[1]{\textcolor[rgb]{0.13,0.47,0.30}{#1}}
\newcommand{\VariableTok}[1]{\textcolor[rgb]{0.07,0.07,0.07}{#1}}
\newcommand{\VerbatimStringTok}[1]{\textcolor[rgb]{0.13,0.47,0.30}{#1}}
\newcommand{\WarningTok}[1]{\textcolor[rgb]{0.37,0.37,0.37}{\textit{#1}}}

\usepackage{longtable,booktabs,array}
\newcounter{none} % for unnumbered tables
\usepackage{calc} % for calculating minipage widths
% Correct order of tables after \paragraph or \subparagraph
\usepackage{etoolbox}
\makeatletter
\patchcmd\longtable{\par}{\if@noskipsec\mbox{}\fi\par}{}{}
\makeatother
% Allow footnotes in longtable head/foot
\IfFileExists{footnotehyper.sty}{\usepackage{footnotehyper}}{\usepackage{footnote}}
\makesavenoteenv{longtable}
\usepackage{graphicx}
\makeatletter
\newsavebox\pandoc@box
\newcommand*\pandocbounded[1]{% scales image to fit in text height/width
  \sbox\pandoc@box{#1}%
  \Gscale@div\@tempa{\textheight}{\dimexpr\ht\pandoc@box+\dp\pandoc@box\relax}%
  \Gscale@div\@tempb{\linewidth}{\wd\pandoc@box}%
  \ifdim\@tempb\p@<\@tempa\p@\let\@tempa\@tempb\fi% select the smaller of both
  \ifdim\@tempa\p@<\p@\scalebox{\@tempa}{\usebox\pandoc@box}%
  \else\usebox{\pandoc@box}%
  \fi%
}
% Set default figure placement to htbp
\def\fps@figure{htbp}
\makeatother



\ifLuaTeX
\usepackage[bidi=basic,shorthands=off]{babel}
\else
\usepackage[bidi=default,shorthands=off]{babel}
\fi
\ifLuaTeX
  \usepackage{selnolig} % disable illegal ligatures
\fi


\setlength{\emergencystretch}{3em} % prevent overfull lines

\providecommand{\tightlist}{%
  \setlength{\itemsep}{0pt}\setlength{\parskip}{0pt}}



 


\KOMAoption{captions}{tableheading}
\makeatletter
\@ifpackageloaded{tcolorbox}{}{\usepackage[skins,breakable]{tcolorbox}}
\@ifpackageloaded{fontawesome5}{}{\usepackage{fontawesome5}}
\definecolor{quarto-callout-color}{HTML}{909090}
\definecolor{quarto-callout-note-color}{HTML}{0758E5}
\definecolor{quarto-callout-important-color}{HTML}{CC1914}
\definecolor{quarto-callout-warning-color}{HTML}{EB9113}
\definecolor{quarto-callout-tip-color}{HTML}{00A047}
\definecolor{quarto-callout-caution-color}{HTML}{FC5300}
\definecolor{quarto-callout-color-frame}{HTML}{acacac}
\definecolor{quarto-callout-note-color-frame}{HTML}{4582ec}
\definecolor{quarto-callout-important-color-frame}{HTML}{d9534f}
\definecolor{quarto-callout-warning-color-frame}{HTML}{f0ad4e}
\definecolor{quarto-callout-tip-color-frame}{HTML}{02b875}
\definecolor{quarto-callout-caution-color-frame}{HTML}{fd7e14}
\makeatother
\makeatletter
\@ifpackageloaded{bookmark}{}{\usepackage{bookmark}}
\makeatother
\makeatletter
\@ifpackageloaded{caption}{}{\usepackage{caption}}
\AtBeginDocument{%
\ifdefined\contentsname
  \renewcommand*\contentsname{Table of contents}
\else
  \newcommand\contentsname{Table of contents}
\fi
\ifdefined\listfigurename
  \renewcommand*\listfigurename{List of Figures}
\else
  \newcommand\listfigurename{List of Figures}
\fi
\ifdefined\listtablename
  \renewcommand*\listtablename{List of Tables}
\else
  \newcommand\listtablename{List of Tables}
\fi
\ifdefined\figurename
  \renewcommand*\figurename{Figure}
\else
  \newcommand\figurename{Figure}
\fi
\ifdefined\tablename
  \renewcommand*\tablename{Table}
\else
  \newcommand\tablename{Table}
\fi
}
\@ifpackageloaded{float}{}{\usepackage{float}}
\floatstyle{ruled}
\@ifundefined{c@chapter}{\newfloat{codelisting}{h}{lop}}{\newfloat{codelisting}{h}{lop}[chapter]}
\floatname{codelisting}{Listing}
\newcommand*\listoflistings{\listof{codelisting}{List of Listings}}
\makeatother
\makeatletter
\makeatother
\makeatletter
\@ifpackageloaded{caption}{}{\usepackage{caption}}
\@ifpackageloaded{subcaption}{}{\usepackage{subcaption}}
\makeatother
\usepackage{bookmark}
\IfFileExists{xurl.sty}{\usepackage{xurl}}{} % add URL line breaks if available
\urlstyle{same}
\hypersetup{
  pdftitle={Predictive Modeling of Residential Property Prices in Pierce County},
  pdfauthor={Edberto Moura Lima; Gerta Hajrullahi; Helga Deutschmann; Mira Radakovic},
  pdflang={en},
  colorlinks=true,
  linkcolor={blue},
  filecolor={Maroon},
  citecolor={Blue},
  urlcolor={Blue},
  pdfcreator={LaTeX via pandoc}}


\title{Predictive Modeling of Residential Property Prices in Pierce
County}
\author{Edberto Moura Lima \and Gerta Hajrullahi \and Helga
Deutschmann \and Mira Radakovic}
\date{2026-05-11}
\begin{document}
\maketitle

\renewcommand*\contentsname{Table of contents}
{
\setcounter{tocdepth}{2}
\tableofcontents
}

\bookmarksetup{startatroot}

\chapter*{Preface}\label{preface}
\addcontentsline{toc}{chapter}{Preface}

\markboth{Preface}{Preface}

This documentation presents the project proposal and development
progress for the Predictive Modeling of Residential Property Prices in
Pierce County, developed as part of the Solution Engineering course at
FHTW.

The project aims to build a robust machine learning pipeline to estimate
housing prices by integrating diverse data sources from the Pierce
County Assessor-Treasurer. This book serves as the central repository
for project objectives, data descriptions, organizational structure, and
development plans.

\section*{Project Team}\label{project-team}
\addcontentsline{toc}{section}{Project Team}

\markright{Project Team}

\begin{itemize}
\tightlist
\item
  \textbf{Mira Radakovic} (Project Leader)
\item
  \textbf{Edberto Moura Lima} (Data Analyst)
\item
  \textbf{Gerta Hajrullahi} (Data Analyst)
\item
  \textbf{Helga Deutschmann} (Data Analyst)
\end{itemize}

\section*{Document Status}\label{document-status}
\addcontentsline{toc}{section}{Document Status}

\markright{Document Status}

This version corresponds to the \textbf{Pr1oject Proposal (Report 1)}.

\begin{itemize}
\tightlist
\item[$\boxtimes$]
  Project Objectives (Qualitative \& Quantitative)
\item[$\boxtimes$]
  Data Source Description \& Dictionary
\item[$\boxtimes$]
  Project Organization \& Task Allocation
\item[$\boxtimes$]
  Development Plan \& Workflow
\item[$\boxtimes$]
  Gantt Schedule \& Milestones
\end{itemize}

\bookmarksetup{startatroot}

\chapter{Project Objective}\label{project-objective}

\section{Background \& Motivation}\label{background-motivation}

Estimating house prices is a classic regression problem with clear
real-world importance. Housing markets are influenced by many
interconnected factors, including property characteristics, location,
neighborhood effects, and market trends over time. Because nearby
properties often share similar values and market conditions change
continuously, house-price prediction provides an excellent example for
applying the complete data-science workflow covered in this course.

For this project, the publicly available Assessor--Treasurer data from
Pierce County is used as the main data source. The dataset contains
detailed and realistic information about properties, transactions, and
assessments across multiple linked tables, offering a large-scale and
practical environment for data exploration, pre-processing, feature
engineering, and predictive modeling.

\section{Project Goal}\label{project-goal}

Develop a machine learning system capable of predicting residential
property prices using structural, spatial, economic, and neighborhood
characteristics.

\subsection{Qualitative Objectives}\label{qualitative-objectives}

\begin{enumerate}
\def\labelenumi{\arabic{enumi}.}
\item
  Build a clean housing dataset for Pierce County.
\item
  Explore which factors influence house prices.
\item
  Engineer predictive variables from property and spatial data.
\item
  Compare traditional statistical and machine learning approaches.
\item
  Deploy a simple Shiny app for interactive price estimation.
\item
  Allow non-technical users to test predictions using custom property
  inputs.
\item
  Create visual outputs that support decision-making.
\end{enumerate}

\section{Quantitative Objectives}\label{quantitative-objectives}

{\def\LTcaptype{none} % do not increment counter
\begin{longtable}[]{@{}
  >{\raggedright\arraybackslash}p{(\linewidth - 4\tabcolsep) * \real{0.4074}}
  >{\raggedright\arraybackslash}p{(\linewidth - 4\tabcolsep) * \real{0.2963}}
  >{\raggedright\arraybackslash}p{(\linewidth - 4\tabcolsep) * \real{0.2963}}@{}}
\toprule\noalign{}
\begin{minipage}[b]{\linewidth}\raggedright
Objective
\end{minipage} & \begin{minipage}[b]{\linewidth}\raggedright
Metric
\end{minipage} & \begin{minipage}[b]{\linewidth}\raggedright
Target
\end{minipage} \\
\midrule\noalign{}
\endhead
\bottomrule\noalign{}
\endlastfoot
\textbf{Model Accuracy} & Root Mean Square Error (RMSE) & \textless{}
USD 75,000 \\
\textbf{Explanatory Power} & Coefficient of Determination (R²) &
\textgreater{} 0.85 \\
\textbf{Model Diversity} & Algorithm Benchmarking & ≥ 4 distinct
algorithms \\
\textbf{Feature Richness} & Engineered Predictors & ≥ 20 model-ready
features \\
\textbf{App Performance} & Real-time Inference & Response time
\textless{} 3 seconds \\
\textbf{Data Integrity} & Table Integration & 100\% of selected source
tables joined \\
\end{longtable}
}

\section{Deliverables \& Timeline}\label{deliverables-timeline}

\begin{itemize}
\tightlist
\item
  \textbf{21.05.2026}: Project Proposal Presentation
\item
  \textbf{06.06.2026}: Interim Results Presentation
\item
  \textbf{24.06.2026}: Final Results Presentation \& Project Report
  Submission
\item
  \textbf{24.06.2026}: Live demo of deployed model App
\end{itemize}

\bookmarksetup{startatroot}

\chapter{Data Source Description \& Data
Dictionary}\label{data-source-description-data-dictionary}

\section{Overview}\label{overview}

The Pierce County Assessor-Treasurer Data Mart contains comprehensive
property assessment and appraisal data maintained by Pierce County,
Washington. The data is organized into multiple related tables that
track property characteristics, improvements, valuations, and sales
transactions.

This data dictionary is a prerequisite for data preparation and serves
as the reference for understanding the relational structure and field
definitions necessary for proper data integration and feature
engineering.

\section{Data Sources}\label{data-sources}

\begin{itemize}
\tightlist
\item
  \textbf{Tax and Assessment System}: Maintains current tax account
  information
\item
  \textbf{Appraisal System}: Maintains detailed property appraisals and
  improvements
\item
  \textbf{Public Records}: Available through Pierce County Public
  Records Request Center
\end{itemize}

\begin{tcolorbox}[enhanced jigsaw, arc=.35mm, bottomrule=.15mm, bottomtitle=1mm, breakable, colback=white, colbacktitle=quarto-callout-note-color!10!white, colframe=quarto-callout-note-color-frame, coltitle=black, left=2mm, leftrule=.75mm, opacityback=0, opacitybacktitle=0.6, rightrule=.15mm, title=\textcolor{quarto-callout-note-color}{\faInfo}\hspace{0.5em}{Note}, titlerule=0mm, toprule=.15mm, toptitle=1mm]

RCW 42.56.070(9) may restrict the release of taxpayer names and mailing
addresses in certain circumstances. Contact the Pierce County Public
Records Request Center for additional information.

\end{tcolorbox}

\begin{center}\rule{0.5\linewidth}{0.5pt}\end{center}

\section{Data Mart Architecture}\label{data-mart-architecture}

\subsection{Table Relationships}\label{table-relationships}

The data mart uses a relational schema with the following structure:

\begin{verbatim}
Tax Account (Parcel Number)
    ↓
Appraisal Account (Parcel Number)
    ├─ Tax Description (Parcel Number, Line Number)
    ├─ Land Attribute (Parcel Number, Attribute)
    ├─ Improvement (Parcel Number, Building ID)
    │   ├─ Improvement Built-As (Parcel Number, Building ID, Built-As Number)
    │   └─ Improvement Detail (Parcel Number, Building ID, Detail Type)
    ├─ Seg/Merge (Seg/Merge Number, Parcel Number)
    └─ Sale (ETN, Parcel Number)

Taxpayer (Parcel Number)
\end{verbatim}

\begin{tcolorbox}[enhanced jigsaw, arc=.35mm, bottomrule=.15mm, bottomtitle=1mm, breakable, colback=white, colbacktitle=quarto-callout-tip-color!10!white, colframe=quarto-callout-tip-color-frame, coltitle=black, left=2mm, leftrule=.75mm, opacityback=0, opacitybacktitle=0.6, rightrule=.15mm, title=\textcolor{quarto-callout-tip-color}{\faLightbulb}\hspace{0.5em}{Key}, titlerule=0mm, toprule=.15mm, toptitle=1mm]

Field names in \textbf{bold} represent primary key(s).

\end{tcolorbox}

\begin{center}\rule{0.5\linewidth}{0.5pt}\end{center}

\section{Table Definitions}\label{table-definitions}

\subsection{Tax Account Table}\label{tax-account-table}

Tracks current tax and assessment account information for each property.

{\def\LTcaptype{none} % do not increment counter
\begin{longtable}[]{@{}
  >{\raggedright\arraybackslash}p{(\linewidth - 6\tabcolsep) * \real{0.2500}}
  >{\raggedright\arraybackslash}p{(\linewidth - 6\tabcolsep) * \real{0.2500}}
  >{\raggedright\arraybackslash}p{(\linewidth - 6\tabcolsep) * \real{0.2500}}
  >{\raggedright\arraybackslash}p{(\linewidth - 6\tabcolsep) * \real{0.2500}}@{}}
\toprule\noalign{}
\begin{minipage}[b]{\linewidth}\raggedright
Field Name
\end{minipage} & \begin{minipage}[b]{\linewidth}\raggedright
Type
\end{minipage} & \begin{minipage}[b]{\linewidth}\raggedright
Size
\end{minipage} & \begin{minipage}[b]{\linewidth}\raggedright
Description
\end{minipage} \\
\midrule\noalign{}
\endhead
\bottomrule\noalign{}
\endlastfoot
\textbf{Parcel Number} & Text & 10 & Unique 10 digit number assigned to
each property \\
Account Type & Text & 5 & REAL - Real Property; STRUC - Structures; PERS
- Personal Property; MOBIL - Mobile Home \\
Property Type & Text & 5 & ASIMP - Administrative Seg Improvement; LNDIM
- Land and Improvements; MBLHM - Mobile Home; RRLEA - Railroad Leases;
SAPP - State Assessed Personal Property; SARP - State Assessed Real
Property; STRUC - Leased Land and/or Structure \\
Site Address & Text & 50 & Physical location address of the property.
May be different than mailing address. ``XXX'' indicates estimated
address using GIS \\
Use Code & Text & 5 & 4-digit numeric code for present highest and best
use of the property \\
Use Description & Text & 50 & Text description of the numeric use
code \\
Tax Year - Prior & Integer & & Tax year for which all prior year values
apply \\
Tax Code Area - Prior Year & Text & 10 & Code describing geographical
area and specific combination of taxing districts. Subject to annual
change \\
Exemption Type - Prior Year & Text & 40 & Type of exemption, if any,
applied to the property \\
Current Use Code - Prior Year & Text & 5 & AGRI - Agricultural; CT -
Current Use Timber; FORDG - Forest Designated; OPBRS - Open Space PBRS;
OPEN - Open Space \\
Land Value - Prior Year & Integer & & Total fair market value of land,
as determined by appraisal process \\
Improvement Value - Prior Year & Integer & & Total fair market value of
improvements (buildings, structures), as determined by appraisal
process \\
Total Market Value - Prior Year & Integer & & Total fair market value of
land and improvements, as determined by appraisal process \\
Taxable Value - Prior Year & Integer & & Market value of the property,
minus any exemptions granted \\
Tax Year - Current & Integer & & Tax year for which all current year
values apply \\
Tax Code Area - Current Year & Text & 10 & Code describing geographical
area and specific combination of taxing districts \\
Exemption Type - Current Year & Text & 40 & Type of exemption, if any,
applied to the property \\
Current Use Code - Current Year & Text & 5 & AGRI - Agricultural; CT -
Current Use Timber; FORDG - Forest Designated; OPBRS - Open Space PBRS;
OPEN - Open Space \\
Land Value - Current Year & Integer & & Total fair market value of land
for the parcel \\
Improvement Value - Current Year & Integer & & Total fair market value
of improvements (buildings, structures) for the parcel \\
Total Market Value - Current Year & Integer & & Total fair market value
of land and improvements for the parcel \\
Taxable Value - Current Year & Integer & & Market value of the property,
minus any exemptions granted \\
Range & Text & 10 & Range \\
Township & Text & 10 & Township \\
Section & Text & 20 & Section \\
Quarter Section & Text & 20 & Quarter \\
Subdivision Name & Text & 50 & Subdivision name \\
Located On Parcel & Text & 10 & Real property parcel number on which
this improvement-only parcel is located (Personal Property, Mobile
Homes, Leasehold Structures) \\
\end{longtable}
}

\begin{center}\rule{0.5\linewidth}{0.5pt}\end{center}

\subsection{Appraisal Account Table}\label{appraisal-account-table}

Detailed property appraisal data with characteristics and valuations.

{\def\LTcaptype{none} % do not increment counter
\begin{longtable}[]{@{}
  >{\raggedright\arraybackslash}p{(\linewidth - 6\tabcolsep) * \real{0.2500}}
  >{\raggedright\arraybackslash}p{(\linewidth - 6\tabcolsep) * \real{0.2500}}
  >{\raggedright\arraybackslash}p{(\linewidth - 6\tabcolsep) * \real{0.2500}}
  >{\raggedright\arraybackslash}p{(\linewidth - 6\tabcolsep) * \real{0.2500}}@{}}
\toprule\noalign{}
\begin{minipage}[b]{\linewidth}\raggedright
Field Name
\end{minipage} & \begin{minipage}[b]{\linewidth}\raggedright
Type
\end{minipage} & \begin{minipage}[b]{\linewidth}\raggedright
Size
\end{minipage} & \begin{minipage}[b]{\linewidth}\raggedright
Description
\end{minipage} \\
\midrule\noalign{}
\endhead
\bottomrule\noalign{}
\endlastfoot
\textbf{Parcel Number} & Text & 10 & Unique 10 digit number assigned to
each property \\
Appraisal Account Type & Text & 15 & Com Condo, Com Leasehold, Com Multi
Unit, Commercial, Industrial, Mobile Home, Reference, Res Com Condo, Res
Leasehold, Residential, Trended Invest \\
Business Name & Text & 50 & Main business name on the parcel, or space
number for mobile/manufactured home in a park \\
Value Area ID & Text & 10 & Identifies the Physical Inspection cycle for
each parcel (PI1 to PI6) \\
Land Economic Area & Text & 30 & Also known as LEA, ties parcels to land
valuation models. For Commercial: identifies land use model by PI year.
For Residential: location code identifying appraisal area, sector and
subsector. For Res Com Condos: relates directly to account land value \\
Buildings & Integer & & Represents the number of improvements
(buildings) located on a parcel \\
Group Account Number & Text & 30 & Identifies and ties all parcels
qualifying as one economic unit (contiguous assessment). Also identifies
all parcels in a condominium project \\
Land Gross Acres & Decimal & 15.4 & Total acres that make up the
economic unit which could include multiple parcels \\
Land Net Acres & Decimal & 15.4 & Size of the individual parcel in
acres \\
Land Gross Square Feet & Decimal & 15.4 & Total square feet that make up
the economic unit which could include multiple parcels \\
Land Net Square Feet & Decimal & 15.4 & Size of the individual parcel in
square feet \\
Land Gross Front Feet & Decimal & 15.4 & For commercial parcels: front
feet of the site facing the arterial. For residential: front feet of the
lot facing the address side. For waterfront: effective waterfront length
of individual lot or combination of lots under single ownership \\
Land Width & Integer & & Residential Use Only. Equal to the Net
Effective Waterfront Front Feet of the individual parcel, or using Group
Account allocation method, then the sum of Net Effective Waterfront
Front Feet for the group \\
Land Depth & Integer & & Residential Use Only. Equal to the Net
Effective Depth of the individual waterfront parcel, or using Group
Account allocation method, then the sum of Net Effective Depths for the
waterfront group. Measured from the top of the bank (where bulkhead
begins) to the inland property line \\
Submerged Area Square Feet & Integer & & Area of a parcel that is
submerged. Not included in the Gross or Net Square Feet \\
Appraisal Date & Date & & Indicates the most recent date an appraiser
physically observed the property and/or made changes to the record \\
Waterfront Type & Text & 30 & Describes the type of waterfront the
property adjoins or has legal access to \\
View Quality & Text & 30 & Assigned to reflect the market appeal of the
overall view available from the dwelling or property \\
Utility Electric & Text & 30 & Identifies if power is installed,
available or is not available on the property \\
Utility Sewer & Text & 30 & Identifies if sewer/septic is installed,
available or not available if the property does not support an on site
sewage disposal system (no perc) \\
Utility Water & Text & 30 & Identifies if water is installed, available
or is not available \\
Street Type & Text & 30 & Identifies if the access street is paved or
unpaved. If blank, the parcel could be a reference parcel or a mobile
home assessed as personal property \\
Latitude & Decimal & 10.5 & The latitude of the parcel centroid. Some
parcels may not have latitude/longitude coordinates (building-only,
mineral rights and reference parcels, as well as others not represented
in the GIS system) \\
Longitude & Decimal & 10.5 & The longitude of the parcel centroid \\
\end{longtable}
}

\begin{center}\rule{0.5\linewidth}{0.5pt}\end{center}

\subsection{Improvement Table}\label{improvement-table}

Detailed building improvement data linked to appraisal accounts.

{\def\LTcaptype{none} % do not increment counter
\begin{longtable}[]{@{}
  >{\raggedright\arraybackslash}p{(\linewidth - 6\tabcolsep) * \real{0.2500}}
  >{\raggedright\arraybackslash}p{(\linewidth - 6\tabcolsep) * \real{0.2500}}
  >{\raggedright\arraybackslash}p{(\linewidth - 6\tabcolsep) * \real{0.2500}}
  >{\raggedright\arraybackslash}p{(\linewidth - 6\tabcolsep) * \real{0.2500}}@{}}
\toprule\noalign{}
\begin{minipage}[b]{\linewidth}\raggedright
Field Name
\end{minipage} & \begin{minipage}[b]{\linewidth}\raggedright
Type
\end{minipage} & \begin{minipage}[b]{\linewidth}\raggedright
Size
\end{minipage} & \begin{minipage}[b]{\linewidth}\raggedright
Description
\end{minipage} \\
\midrule\noalign{}
\endhead
\bottomrule\noalign{}
\endlastfoot
\textbf{Parcel Number} & Text & 10 & Unique 10 digit number assigned to
each property \\
\textbf{Building ID} & Text & 5 & Unique identifier for each improvement
record on a parcel. Not all improvements are on a separate record and
may not always be in sequence \\
Property Type & Text & 15 & Links buildings to proper cost and
depreciation tables. Options include: Commercial, Duplex, Industrial,
Mobile Home, Multiple Unit, Out Building, Residential, Townhouse,
Triplex \\
Neighborhood & Text & 10 & Represents boundaries for land subject to
similar social, environmental, economic and governmental forces.
Boundaries are typically located within a geographic area defined by
natural, man-made, or political boundaries. Residential Neighborhoods
and LEA's are the same. Commercial Neighborhoods define the location of
the property while Commercial LEA's define the use of the property \\
Neighborhood Extension & Text & 10 & Based on property type: Commercial
- represents the primary occupancy or use of the property. Residential -
not used, therefore they are set to zero. Mobile/Manufactured Homes -
based upon the underlying land ownership. All Mobile/Manufactured homes
located on land owned by the Mobile/Manufactured home owner are coded
with the Neighborhood Extension 0. Exception - Mobile/Manufactured homes
located within a park and owned by the park owner, will carry the same
Neighborhood Extension as the park. Mobile Homes in a park have a
neighborhood extension of: PL 1 (P = Park, L = Leased); P2, P3, P4,
etc \\
Square Feet & Integer & & Sum of the square feet for the `built as'
types identified for the building. Each building may have one or more
`built as' type (see IMPROVEMENT BUILTAS table), though usually only
one. The exceptions are `built as' \# 124 and 126 (Add on only Res \&
Com) or any type of Storage Tank, where it represents the count of those
`built as' as opposed to the sum of their square feet \\
Net Square Feet & Integer & & Used primarily to calculate income on
commercial properties, where it typically represents the sum of the net
rentable area(s). This number may be representative of one or more
buildings. If part of a Group Account, these buildings may be located on
more than one parcel. For Residential it is equal to the built as square
feet or zero (0) \\
\end{longtable}
}

\begin{center}\rule{0.5\linewidth}{0.5pt}\end{center}

\subsection{Improvement Built-As
Table}\label{improvement-built-as-table}

Detailed specifications for building construction styles and features.

{\def\LTcaptype{none} % do not increment counter
\begin{longtable}[]{@{}
  >{\raggedright\arraybackslash}p{(\linewidth - 6\tabcolsep) * \real{0.2500}}
  >{\raggedright\arraybackslash}p{(\linewidth - 6\tabcolsep) * \real{0.2500}}
  >{\raggedright\arraybackslash}p{(\linewidth - 6\tabcolsep) * \real{0.2500}}
  >{\raggedright\arraybackslash}p{(\linewidth - 6\tabcolsep) * \real{0.2500}}@{}}
\toprule\noalign{}
\begin{minipage}[b]{\linewidth}\raggedright
Field Name
\end{minipage} & \begin{minipage}[b]{\linewidth}\raggedright
Type
\end{minipage} & \begin{minipage}[b]{\linewidth}\raggedright
Size
\end{minipage} & \begin{minipage}[b]{\linewidth}\raggedright
Description
\end{minipage} \\
\midrule\noalign{}
\endhead
\bottomrule\noalign{}
\endlastfoot
\textbf{Parcel Number} & Text & 10 & Unique 10 digit number assigned to
each property \\
\textbf{Building ID} & Text & 5 & Unique identifier for each improvement
record on a parcel \\
\textbf{Built-As Number} & Integer & & Unique number identifying each
Built-As on an improvement. Sequential number assigned by the system
starting with the largest Built-As being number 1 \\
Built-As ID & Integer & & Number associated with the description of the
purpose or style of construction of the building \\
Built-As Description & Text & 255 & Original purpose for and/or style of
construction of the building. The Built As is associated with the
Marshall and Swift cost and depreciation tables \\
Built-As Square Feet & Integer & & Sum of the square feet for the `built
as' types identified for the building. Each building may have one or
more `built as' type, though usually only one. The exceptions are `built
as' \# 124 and 126 (Add on only Res \& Com) or any type of Storage
Tank \\
HVAC & Integer & & Code associated with the predominant heating source
for the built-as structure \\
HVAC Description & Text & 30 & Text description associated with the
predominant heating source for the built-as structure (Forced Air,
Electric Baseboard, Steam, etc) \\
Exterior & Text & 25 & Predominant type of construction materials used
for the exterior siding on Residential Buildings \\
Interior & Text & 15 & Predominant type of materials used on the
interior walls (Sheetrock or Paneling). Collected for Residential and
Mobile Homes only \\
Stories & Decimal & 13.2 & Number of floors/building levels above grade.
Stories do not include attic or basement areas \\
Story Height & Integer & & Based on property type: Commercial -
determined by taking an average of the overall ceiling heights of all
floors for the Blt As. Residential - determined by the majority of the
story height for the first floor (main). Story height is automatically
rounded by the system (The system rounds up for .5 and higher and rounds
down for anything below .5) \\
Sprinkler Square Feet & Integer & & Total square footage covered by a
sprinkler system. Sprinkler SF does not always equal the built as square
footage \\
Roof Cover & Text & 20 & Material used for the roof (Composition
Shingles, Wood Shake, Concrete Tile, etc) \\
Bedrooms & Integer & & Number of bedrooms listed for a residential
property (Collected for informational purposes only) \\
Bathrooms & Decimal & 7.2 & Number of baths listed for a residential
property. The number is listed as a decimal, i.e.~2.75 = two full and
one three-quarter baths. A tub/sink/toilet combination (plus any
additional fixtures) is considered 1.0 bath. A shower/sink/toilet
combination (plus any additional fixtures) is 0.75 bath. A sink/toilet
combination is .5 bath \\
Units & Integer & & Number of separate units within the building. For
Commercial properties, Units may be combined to generate an income
approach for the economic unit on building one \\
Class Code & Text & 20 & Code for one of five basic cost groups by type
of framing (supporting beams and columns), walls, floors and roof
structures, and fireproofing. Typically Commercial Buildings \\
Class Description & Text & 50 & Text description for one of five basic
cost groups by type of framing (supporting beams and columns), walls,
floors and roof structures, and fireproofing. Typically Commercial
Buildings \\
Year Built & Integer & & Year the building was built, as stated by the
building permit or a historical record \\
Year Remodeled & Integer & & Year in which improvements are made to the
existing structure that extend the life of the structure \\
Adjusted Year Built & Integer & & If greater than the Year Built, it
indicates improvements have been made to the original property over and
above normal maintenance which would effectively reduce the age of the
building \\
Physical Age & Integer & & Tax year less the adjusted year built equals
the physical age of the built as \\
Built-As Length & Integer & & Length of the mobile/manufactured home \\
Built-As Width & Integer & & Width of the mobile/manufactured home \\
Mobile Home Model & Text & 50 & Model of the mobile/manufactured home \\
\end{longtable}
}

\begin{center}\rule{0.5\linewidth}{0.5pt}\end{center}

\subsection{Improvement Detail Table}\label{improvement-detail-table}

Additional specific improvement characteristics and features.

{\def\LTcaptype{none} % do not increment counter
\begin{longtable}[]{@{}
  >{\raggedright\arraybackslash}p{(\linewidth - 6\tabcolsep) * \real{0.2500}}
  >{\raggedright\arraybackslash}p{(\linewidth - 6\tabcolsep) * \real{0.2500}}
  >{\raggedright\arraybackslash}p{(\linewidth - 6\tabcolsep) * \real{0.2500}}
  >{\raggedright\arraybackslash}p{(\linewidth - 6\tabcolsep) * \real{0.2500}}@{}}
\toprule\noalign{}
\begin{minipage}[b]{\linewidth}\raggedright
Field Name
\end{minipage} & \begin{minipage}[b]{\linewidth}\raggedright
Type
\end{minipage} & \begin{minipage}[b]{\linewidth}\raggedright
Size
\end{minipage} & \begin{minipage}[b]{\linewidth}\raggedright
Description
\end{minipage} \\
\midrule\noalign{}
\endhead
\bottomrule\noalign{}
\endlastfoot
\textbf{Parcel Number} & Text & 10 & Unique 10 digit number assigned to
each property \\
\textbf{Building ID} & Text & 5 & Unique identifier for each improvement
record on a parcel \\
\textbf{Detail Type} & Text & 10 & Category that specific improvement
characteristics are coded under. Options include: Add On, Appliance,
Balcony, Basement, Carport, Elevator, Fixture, Garage, Mezzanine, Porch,
Rough In, Storage. The Add On type allows the appraiser to assign a
quality different from the main structure for the characteristics within
this field \\
Detail Description & Text & 50 & Brief explanation of the improvement
characteristic (Fireplace, PreFab/Stoves, Laundry Facility, Type of
Porch (Cvrd Wd Deck), Type of Bath (Bath 2 Fixture)) \\
Units & Decimal & 15.4 & Component parts of the building
characteristics. The units indicate either the square footage of the
structure or a count of each component. Examples: Appliance, Allowance,
1 (count) or Porch, Open Slab, 337 (sf), or Fixture -- Bath 3 Fixtures
-- 2 (count) \\
\end{longtable}
}

\begin{center}\rule{0.5\linewidth}{0.5pt}\end{center}

\subsection{Land Attribute Table}\label{land-attribute-table}

Land-specific attributes and characteristics affecting valuation.

{\def\LTcaptype{none} % do not increment counter
\begin{longtable}[]{@{}
  >{\raggedright\arraybackslash}p{(\linewidth - 6\tabcolsep) * \real{0.2500}}
  >{\raggedright\arraybackslash}p{(\linewidth - 6\tabcolsep) * \real{0.2500}}
  >{\raggedright\arraybackslash}p{(\linewidth - 6\tabcolsep) * \real{0.2500}}
  >{\raggedright\arraybackslash}p{(\linewidth - 6\tabcolsep) * \real{0.2500}}@{}}
\toprule\noalign{}
\begin{minipage}[b]{\linewidth}\raggedright
Field Name
\end{minipage} & \begin{minipage}[b]{\linewidth}\raggedright
Type
\end{minipage} & \begin{minipage}[b]{\linewidth}\raggedright
Size
\end{minipage} & \begin{minipage}[b]{\linewidth}\raggedright
Description
\end{minipage} \\
\midrule\noalign{}
\endhead
\bottomrule\noalign{}
\endlastfoot
\textbf{Parcel Number} & Text & 10 & Unique 10 digit number assigned to
each property \\
\textbf{Attribute} & Text & 30 & Describes the category that a land
attribute is assigned to: Residential (R), Commercial (C) or Res Com
Condo (X) and the subcategory, such as Amenities, Economic, Functional,
Neighborhood, Utilities, etc \\
Attribute Description & Text & 30 & Further defines the land attribute.
Land attributes may or may not contribute to value (Examples of
attributes: View Avg, NBHD Gated, WF Bank Low, water available etc) \\
\end{longtable}
}

\begin{center}\rule{0.5\linewidth}{0.5pt}\end{center}

\subsection{Sale Table}\label{sale-table}

Historical property sales transaction data.

{\def\LTcaptype{none} % do not increment counter
\begin{longtable}[]{@{}
  >{\raggedright\arraybackslash}p{(\linewidth - 6\tabcolsep) * \real{0.2500}}
  >{\raggedright\arraybackslash}p{(\linewidth - 6\tabcolsep) * \real{0.2500}}
  >{\raggedright\arraybackslash}p{(\linewidth - 6\tabcolsep) * \real{0.2500}}
  >{\raggedright\arraybackslash}p{(\linewidth - 6\tabcolsep) * \real{0.2500}}@{}}
\toprule\noalign{}
\begin{minipage}[b]{\linewidth}\raggedright
Field Name
\end{minipage} & \begin{minipage}[b]{\linewidth}\raggedright
Type
\end{minipage} & \begin{minipage}[b]{\linewidth}\raggedright
Size
\end{minipage} & \begin{minipage}[b]{\linewidth}\raggedright
Description
\end{minipage} \\
\midrule\noalign{}
\endhead
\bottomrule\noalign{}
\endlastfoot
\textbf{ETN} & Text & 11 & Recording number issued by the Auditor's
Office when a property transaction is recorded. Excise tax must be paid
before the document of conveyance can be recorded or the mobile home
title transferred \\
Parcel Count & Integer & & Number of parcels associated with a sale \\
Parcel Number & Text & 10 & Unique 10 digit number assigned to each
property \\
Sale Date & Date & & Date the legal document (deed) was executed \\
Sale Price & Decimal & 15.2 & Dollar amount recorded on the ETN \\
Deed Type & Text & 40 & Type of document which conveyed an interest in
or legal title to the property \\
Grantor & Text & 80 & Individual(s) or company conveying ownership or
interest in the property described on the deed \\
Grantee & Text & 80 & Individual(s) or company purchasing ownership or
interest in the property described on the deed \\
Valid/Invalid & Text & 7 & Used to code the validity of a sale for
appraisal purposes. A sale in the open market between two unrelated
parties, each of whom is reasonably knowledgeable of market conditions
and under no undue pressure to buy or sell would be considered a valid
transaction. However, a sale can be a valid transaction but if the
conditions of the sale fit one of the 27 deletion categories identified
in the State Ratio RCW, the sale would be coded invalid for Assessor
purposes (sale between relatives, estate sale, percent interest, etc) \\
Confirmed/Unconfirmed & Text & 11 & Identifies if the property
characteristics at the time of the sale and/or circumstances of the sale
were verified by Assessor-Treasurer staff \\
Exclude Reason & Text & 30 & State required description of why a sales
transaction must be excluded from sales studies. Examples include
non-arms length sales transactions such as ``Family different last
names'', sale transactions with undue pressure such as ``Estate sale''
and ``Forced Sale Trans in Lieu Fcl'', and sales where the property
characteristics have changed since the sale, ``Improved after sale'' \\
Improved/Vacant & Text & 8 & If any parcel in the sale has any building
value it is coded as improved. Vacant is vacant land with no building
value \\
Appraisal Account Type & Text & 15 & How a parcel is classified (Com
Condo, Com Leasehold, Com Multi Unit, Commercial, Industrial, Mobile
Home, Reference, Res Com Condo, Res Leasehold, Residential, Trended
Invest) \\
\end{longtable}
}

\begin{center}\rule{0.5\linewidth}{0.5pt}\end{center}

\subsection{Seg/Merge Table}\label{segmerge-table}

Property segregation and merger transaction history.

{\def\LTcaptype{none} % do not increment counter
\begin{longtable}[]{@{}
  >{\raggedright\arraybackslash}p{(\linewidth - 6\tabcolsep) * \real{0.2500}}
  >{\raggedright\arraybackslash}p{(\linewidth - 6\tabcolsep) * \real{0.2500}}
  >{\raggedright\arraybackslash}p{(\linewidth - 6\tabcolsep) * \real{0.2500}}
  >{\raggedright\arraybackslash}p{(\linewidth - 6\tabcolsep) * \real{0.2500}}@{}}
\toprule\noalign{}
\begin{minipage}[b]{\linewidth}\raggedright
Field Name
\end{minipage} & \begin{minipage}[b]{\linewidth}\raggedright
Type
\end{minipage} & \begin{minipage}[b]{\linewidth}\raggedright
Size
\end{minipage} & \begin{minipage}[b]{\linewidth}\raggedright
Description
\end{minipage} \\
\midrule\noalign{}
\endhead
\bottomrule\noalign{}
\endlastfoot
\textbf{Seg/Merge Number} & Text & 12 & Unique number assigned to each
segregation \\
Parent/Child Indicator & Text & 1 & P - Parent Parcel - inactivated
parcel being replaced by the new child parcel. There may be zero, one or
more parent parcels involved with each segregation. In some cases, a
parent parcel may remain active after the segregation. C - Child Parcel
- a new parcel created as a result of the seg. There are one or more
child(ren) parcels involved with each segregation \\
\textbf{Parcel Number} & Text & 10 & Unique 10 digit number assigned to
each property \\
Continued Indicator & Text & 1 & Y - the parcel is still active after
the segregation; N - the parcel is no longer active after the
segregation \\
Completed Date & Date & & Date that the segregation was completed \\
Tax Year & Text & 4 & Tax year of the segregation \\
\end{longtable}
}

\begin{center}\rule{0.5\linewidth}{0.5pt}\end{center}

\subsection{Taxpayer Table}\label{taxpayer-table}

Property owner information (subject to privacy restrictions).

{\def\LTcaptype{none} % do not increment counter
\begin{longtable}[]{@{}
  >{\raggedright\arraybackslash}p{(\linewidth - 6\tabcolsep) * \real{0.2500}}
  >{\raggedright\arraybackslash}p{(\linewidth - 6\tabcolsep) * \real{0.2500}}
  >{\raggedright\arraybackslash}p{(\linewidth - 6\tabcolsep) * \real{0.2500}}
  >{\raggedright\arraybackslash}p{(\linewidth - 6\tabcolsep) * \real{0.2500}}@{}}
\toprule\noalign{}
\begin{minipage}[b]{\linewidth}\raggedright
Field Name
\end{minipage} & \begin{minipage}[b]{\linewidth}\raggedright
Type
\end{minipage} & \begin{minipage}[b]{\linewidth}\raggedright
Size
\end{minipage} & \begin{minipage}[b]{\linewidth}\raggedright
Description
\end{minipage} \\
\midrule\noalign{}
\endhead
\bottomrule\noalign{}
\endlastfoot
\textbf{Parcel Number} & Text & 10 & Unique 10 digit number assigned to
each property \\
\end{longtable}
}

\begin{tcolorbox}[enhanced jigsaw, arc=.35mm, bottomrule=.15mm, bottomtitle=1mm, breakable, colback=white, colbacktitle=quarto-callout-warning-color!10!white, colframe=quarto-callout-warning-color-frame, coltitle=black, left=2mm, leftrule=.75mm, opacityback=0, opacitybacktitle=0.6, rightrule=.15mm, title=\textcolor{quarto-callout-warning-color}{\faExclamationTriangle}\hspace{0.5em}{Warning}, titlerule=0mm, toprule=.15mm, toptitle=1mm]

RCW 42.56.070(9) may restrict the release of taxpayer names and mailing
addresses in certain circumstances. Please contact the Pierce County
Public Records Request Center to obtain additional information and/or
initiate a request for the taxpayer table.

\end{tcolorbox}

\begin{center}\rule{0.5\linewidth}{0.5pt}\end{center}

\subsection{Tax Description Table}\label{tax-description-table}

Detailed tax coding and descriptions for parcels.

{\def\LTcaptype{none} % do not increment counter
\begin{longtable}[]{@{}
  >{\raggedright\arraybackslash}p{(\linewidth - 6\tabcolsep) * \real{0.2500}}
  >{\raggedright\arraybackslash}p{(\linewidth - 6\tabcolsep) * \real{0.2500}}
  >{\raggedright\arraybackslash}p{(\linewidth - 6\tabcolsep) * \real{0.2500}}
  >{\raggedright\arraybackslash}p{(\linewidth - 6\tabcolsep) * \real{0.2500}}@{}}
\toprule\noalign{}
\begin{minipage}[b]{\linewidth}\raggedright
Field Name
\end{minipage} & \begin{minipage}[b]{\linewidth}\raggedright
Type
\end{minipage} & \begin{minipage}[b]{\linewidth}\raggedright
Size
\end{minipage} & \begin{minipage}[b]{\linewidth}\raggedright
Description
\end{minipage} \\
\midrule\noalign{}
\endhead
\bottomrule\noalign{}
\endlastfoot
\textbf{Parcel Number} & Text & 10 & Unique 10 digit number assigned to
each property \\
\textbf{Line Number} & Integer & & Sequential number identifying each
line of a tax description \\
Tax Description Line & Text & 200 & Tax description text \\
\end{longtable}
}

\begin{center}\rule{0.5\linewidth}{0.5pt}\end{center}

\section{Key Relationships}\label{key-relationships}

\begin{itemize}
\tightlist
\item
  \textbf{One-to-Many}: Each parcel (Parcel Number) can have multiple
  improvements (Building IDs)
\item
  \textbf{One-to-Many}: Each improvement can have multiple built-as
  specifications and detail records
\item
  \textbf{One-to-One}: Each parcel has one appraisal account and one tax
  account
\item
  \textbf{One-to-Many}: Each parcel can have multiple sales transactions
  (ETN)
\item
  \textbf{One-to-Many}: Each parcel can have multiple land attributes
\item
  \textbf{One-to-Many}: Each parcel can be involved in multiple
  seg/merge transactions
\end{itemize}

\begin{center}\rule{0.5\linewidth}{0.5pt}\end{center}

\section{Data Quality Notes}\label{data-quality-notes}

\begin{itemize}
\tightlist
\item
  \textbf{Historical Data}: Seg/Merge and Sale tables contain historical
  data. All other tables contain data for active parcels only
\item
  \textbf{Missing Coordinates}: Some parcels (building-only, mineral
  rights, reference parcels) may not have latitude/longitude coordinates
\item
  \textbf{Address Estimation}: If address contains ``XXX'', it was
  estimated using GIS
\item
  \textbf{Appraisal Dates}: Fields like Appraisal Date indicate when the
  property was last physically inspected
\item
  \textbf{Value Dates}: Land, Improvement, and Total Market Values are
  tracked for both Prior Year and Current Year
\end{itemize}

\begin{center}\rule{0.5\linewidth}{0.5pt}\end{center}

\section{Data Dictionary Notes}\label{data-dictionary-notes}

\begin{itemize}
\tightlist
\item
  \textbf{Field names in bold} represent primary key(s) in their
  respective tables
\item
  \textbf{Size column} represents maximum field length for Text and
  Decimal fields
\item
  Fields without size notation are fixed or variable based on actual
  values
\item
  Most text fields are case-sensitive as recorded in the assessor system
\end{itemize}

\begin{center}\rule{0.5\linewidth}{0.5pt}\end{center}

\section{Access \& Contact}\label{access-contact}

For additional information, data downloads, or to request specific
datasets:

\begin{itemize}
\tightlist
\item
  \textbf{Source}: Pierce County Assessor-Treasurer Office
\item
  \textbf{Website}: \url{https://www.piercecountywa.gov/}
\item
  \textbf{Records Request}: Pierce County Public Records Request Center
\end{itemize}

\bookmarksetup{startatroot}

\chapter{Project Organization \& Task
Allocation}\label{project-organization-task-allocation}

\section{Group Structure}\label{group-structure}

{\def\LTcaptype{none} % do not increment counter
\begin{longtable}[]{@{}ll@{}}
\toprule\noalign{}
Name & Role \\
\midrule\noalign{}
\endhead
\bottomrule\noalign{}
\endlastfoot
\textbf{Mira Radakovic (PL)} & Project Leader \\
\textbf{Edberto Moura Lima} & Data Analyst \\
\textbf{Gerta Hajrullahi} & Data Analyst \\
\textbf{Helga Deutschmann} & Data Analyst \\
\end{longtable}
}

\section{Tools \& Infrastructure}\label{tools-infrastructure}

{\def\LTcaptype{none} % do not increment counter
\begin{longtable}[]{@{}ll@{}}
\toprule\noalign{}
Component & Detail \\
\midrule\noalign{}
\endhead
\bottomrule\noalign{}
\endlastfoot
\textbf{Language} & R (primary) \\
\textbf{Database} & Supabase (PostgreSQL) \\
\textbf{Version Control} & Git + GitHub \\
\textbf{Project Management} & GitHub Projects (Kanban Board) \\
\textbf{Communication} & Group messaging + sync meetings \\
\textbf{Documentation} & Quarto \\
\textbf{Data Wrangling} & dplyr, tidyr, readr \\
\textbf{EDA \& Viz} & ggplot2, corrplot, leaflet, plotly \\
\textbf{Modelling} & caret / tidymodels, ranger, xgboost, glmnet \\
\textbf{Spatial Analysis} & sf, geosphere \\
\textbf{Deployment} & Shiny, shinydashboard \\
\end{longtable}
}

\section{Project Kanban Board}\label{project-kanban-board}

The following board tracks the real-time status of all project tasks.

\pandocbounded{\includegraphics[keepaspectratio]{tasks.png}}

\section{Risk Management}\label{risk-management}

{\def\LTcaptype{none} % do not increment counter
\begin{longtable}[]{@{}
  >{\raggedright\arraybackslash}p{(\linewidth - 2\tabcolsep) * \real{0.5000}}
  >{\raggedright\arraybackslash}p{(\linewidth - 2\tabcolsep) * \real{0.5000}}@{}}
\toprule\noalign{}
\begin{minipage}[b]{\linewidth}\raggedright
Risk
\end{minipage} & \begin{minipage}[b]{\linewidth}\raggedright
Mitigation
\end{minipage} \\
\midrule\noalign{}
\endhead
\bottomrule\noalign{}
\endlastfoot
Large number of tables / complex joins & Build join pipeline
incrementally; validate row counts at each step \\
Many missing values in building details & Imputation strategies
(median/mode); sensitivity analysis on imputed variables \\
Data covers multiple property types (commercial, mobile homes) & Filter
to Residential Improved sales only at ingestion \\
Temporal drift in prices (market cycles) & Include sale year/month as
features; optionally restrict to recent years \\
Skewed Sale Price distribution & Log-transform target; evaluate models
on both scales \\
Shiny App complexity & Start with minimal viable app; add features
iteratively \\
\end{longtable}
}

\bookmarksetup{startatroot}

\chapter{Development Plan Overview}\label{development-plan-overview}

The project is splinted into work packages (WP0--WP7) over 6 weeks.

\textbf{Detailed work package specifications} (tasks, deliverables,
quality criteria, and dependencies) are provided in
\href{05-project-schedule-and-workpackages.qmd}{Chapter 5: Project
Schedule \& Work Packages}.

\section{Workflow Diagram}\label{workflow-diagram}

\begin{figure}

\centering{

\pandocbounded{\includegraphics[keepaspectratio]{workflow.png}}

}

\caption{\label{fig-workflow}Project Development Workflow}

\end{figure}%

\section{Work Package Summary}\label{work-package-summary}

{\def\LTcaptype{none} % do not increment counter
\begin{longtable}[]{@{}lll@{}}
\toprule\noalign{}
WP & Phase & Duration \\
\midrule\noalign{}
\endhead
\bottomrule\noalign{}
\endlastfoot
WP0 & Project Setup & 13--19 May \\
WP1 & Data Integration & 20--26 May \\
WP2 & Data Cleaning & 27 May -- 2 June \\
WP3 & EDA \& Features & 30 May -- 9 June \\
WP4 & Preprocessing & 7--9 June \\
WP5 & Model Development & 9--15 June \\
WP6 & Model Evaluation & 16--20 June \\
WP7 & Deployment & 20--24 June \\
\end{longtable}
}

\begin{center}\rule{0.5\linewidth}{0.5pt}\end{center}

\section{Communication \& Oversight}\label{communication-oversight}

\textbf{Weekly Sync}: - Review completed tasks - Discuss blockers \&
dependencies - Update Kanban boar - Plan next week

\textbf{Status Reporting}: Mira provides weekly status to lecturer

\begin{center}\rule{0.5\linewidth}{0.5pt}\end{center}

\bookmarksetup{startatroot}

\chapter{Project Schedule \& Work
Packages}\label{project-schedule-work-packages}

\textbf{Project Duration}: 13 May -- 24 June 2026 (6 weeks)

\begin{center}\rule{0.5\linewidth}{0.5pt}\end{center}

\section{Interactive Gantt Chart}\label{interactive-gantt-chart}

The following chart outlines the project timeline, milestones, and
critical path.

\includegraphics[width=7.69in,height=3.79in]{05-project-schedule-and-workpackages_files/figure-latex/mermaid-figure-1.png}

\begin{figure}

\centering{

\pandocbounded{\includegraphics[keepaspectratio]{gantt.png}}

}

\caption{\label{fig-gantt-static}Project Gantt Chart (Static Version)}

\end{figure}%

\begin{center}\rule{0.5\linewidth}{0.5pt}\end{center}

\section{Detailed Work Package
Definitions}\label{detailed-work-package-definitions}

\subsection{WP0: Project Setup}\label{wp0-project-setup}

\textbf{Duration}: 13--19 May \textbar{} \textbf{Status}: Done

\subsubsection{Tasks}\label{tasks}

\begin{itemize}
\tightlist
\item
  Define team roles and responsibilities
\item
  Create Git repository (private, shared with lecturer)
\item
  Set up Supabase database connection
\item
  Initialize Kanban board (Backlog/To Do/Doing/Testing/Done)
\item
  Configure R development environment
\end{itemize}

\subsubsection{Deliverables}\label{deliverables}

✓ GitHub repository ✓ Supabase connection ✓ Kanban board ✓ Team charter
document

\begin{center}\rule{0.5\linewidth}{0.5pt}\end{center}

\subsection{WP1: Data Integration}\label{wp1-data-integration}

\textbf{Duration}: 20--26 May \textbar{} \textbf{Status}: To Do

\subsubsection{Tasks}\label{tasks-1}

\begin{enumerate}
\def\labelenumi{\arabic{enumi}.}
\tightlist
\item
  Download Pierce County data files
\item
  Map table relationships and join keys
\item
  Perform sequential table joins (validate at each step)
\item
  Validate referential integrity and row counts
\item
  Create data integration report
\end{enumerate}

\subsubsection{Deliverables}\label{deliverables-1}

\begin{itemize}
\tightlist
\item
  Integrated dataset (.RData)
\item
  Table mapping diagram
\item
  Data dictionary (updated)
\item
  Validation report with row counts
\end{itemize}

\subsubsection{Quality Criteria}\label{quality-criteria}

\begin{itemize}
\tightlist
\item[$\square$]
  All tables joined successfully
\item[$\square$]
  Row counts documented at each step
\item[$\square$]
  No unexpected NULL values
\item[$\square$]
  Referential integrity verified
\end{itemize}

\begin{center}\rule{0.5\linewidth}{0.5pt}\end{center}

\subsection{WP2: Data Cleaning \&
Preparation}\label{wp2-data-cleaning-preparation}

\textbf{Duration}: 27 May -- 2 June \textbar{} \textbf{Status}: To Do

\subsubsection{Tasks}\label{tasks-2}

\begin{enumerate}
\def\labelenumi{\arabic{enumi}.}
\tightlist
\item
  Filter to Residential Improved properties only
\item
  Remove duplicate records
\item
  Handle missing values (document strategy)
\item
  Examine and document outliers
\item
  Variable selection --- keep only analysis-ready columns
\item
  Run validation test suite
\end{enumerate}

\subsubsection{Deliverables}\label{deliverables-2}

\begin{itemize}
\tightlist
\item
  Cleaned dataset (.RData)
\item
  Data cleaning log
\item
  Missing value analysis
\item
  Quality assurance report
\end{itemize}

\subsubsection{Quality Criteria}\label{quality-criteria-1}

\begin{itemize}
\tightlist
\item[$\square$]
  No critical missing values
\item[$\square$]
  Data types verified
\item[$\square$]
  Outliers documented
\item[$\square$]
  Test suite passes
\end{itemize}

\begin{center}\rule{0.5\linewidth}{0.5pt}\end{center}

\subsection{WP3: Exploratory Data Analysis \& Feature
Engineering}\label{wp3-exploratory-data-analysis-feature-engineering}

\textbf{Duration}: 30 May -- 9 June \textbar{} \textbf{Status}: To Do

\subsubsection{WP3a: EDA Tasks}\label{wp3a-eda-tasks}

\begin{itemize}
\tightlist
\item
  Univariate analysis (distributions, summaries)
\item
  Bivariate analysis (price vs.~key variables)
\item
  Spatial analysis with leaflet maps
\item
  Correlation matrix and hypothesis testing
\end{itemize}

\subsubsection{WP3b: Feature Engineering
Tasks}\label{wp3b-feature-engineering-tasks}

\begin{itemize}
\tightlist
\item
  Create derived features (age, price/sqft, neighborhood effects, etc.)
\item
  Encode categorical variables
\item
  Feature selection --- target ≥15 final features
\item
  Document feature importance baseline
\end{itemize}

\subsubsection{Deliverables}\label{deliverables-3}

\begin{itemize}
\tightlist
\item
  EDA report with visualizations
\item
  Correlation heatmap
\item
  Feature engineering log
\item
  Final feature set documentation (≥15 features)
\end{itemize}

\subsubsection{Quality Criteria}\label{quality-criteria-2}

\begin{itemize}
\tightlist
\item[$\square$]
  All visualizations publication-ready
\item[$\square$]
  ≥15 features engineered and named
\item[$\square$]
  No data leakage from target variable
\item[$\square$]
  Feature-target correlations documented
\end{itemize}

\begin{center}\rule{0.5\linewidth}{0.5pt}\end{center}

\subsection{WP4: Preprocessing}\label{wp4-preprocessing}

\textbf{Duration}: 7--9 June \textbar{} \textbf{Status}: To Do

\subsubsection{Tasks}\label{tasks-3}

\begin{enumerate}
\def\labelenumi{\arabic{enumi}.}
\tightlist
\item
  Encode categorical variables (finalize encoding scheme)
\item
  Scale numerical variables (standardize or normalize)
\item
  Train/validation/test split (70/15/15\% stratified)
\item
  Document random seed for reproducibility
\end{enumerate}

\subsubsection{Deliverables}\label{deliverables-4}

\begin{itemize}
\tightlist
\item
  3 split datasets (train, validation, test)
\item
  Encoding mapping tables
\item
  Scaling parameters (saved as R object)
\item
  Data split documentation
\end{itemize}

\subsubsection{Quality Criteria}\label{quality-criteria-3}

\begin{itemize}
\tightlist
\item[$\square$]
  No data leakage (scaling only from training set)
\item[$\square$]
  Balanced distribution across splits
\item[$\square$]
  Random seed documented
\item[$\square$]
  All features properly encoded/scaled
\end{itemize}

\begin{center}\rule{0.5\linewidth}{0.5pt}\end{center}

\subsection{WP5: Model Development (Parallel
Tracks)}\label{wp5-model-development-parallel-tracks}

\textbf{Duration}: 9--15 June \textbar{} \textbf{Status}: To Do

\subsubsection{WP5a: Baseline \& Tree
Models}\label{wp5a-baseline-tree-models}

\begin{itemize}
\tightlist
\item
  Baseline Linear Regression
\item
  Decision Tree (hyperparameter tuning)
\item
  Random Forest (grid search over parameters)
\end{itemize}

\subsubsection{WP5b: Support Vector
Machine}\label{wp5b-support-vector-machine}

\begin{itemize}
\tightlist
\item
  SVM setup (kernel selection)
\item
  Hyperparameter tuning (C, epsilon)
\item
  Cross-validation
\end{itemize}

\subsubsection{WP5c: k-Nearest
Neighbours}\label{wp5c-k-nearest-neighbours}

\begin{itemize}
\tightlist
\item
  Setup (distance metric, k selection)
\item
  kNN based on: distance (lon, lat) + square metre
\item
  Hyperparameter optimization
\end{itemize}

\subsubsection{WP5d: Linear Regression}\label{wp5d-linear-regression}

\begin{itemize}
\tightlist
\item
  Full feature set training
\item
  Feature importance extraction
\item
  Residual diagnostics
\end{itemize}

\subsubsection{Deliverables}\label{deliverables-5}

\begin{itemize}
\tightlist
\item
  Trained model objects (.RDS)
\item
  Cross-validation results
\item
  Validation \& test metrics
\item
  Model configuration documentation
\end{itemize}

\subsubsection{Quality Criteria}\label{quality-criteria-4}

\begin{itemize}
\tightlist
\item[$\square$]
  All 5 models trained and saved
\item[$\square$]
  Cross-validation RMSE \textless{} 100k (preliminary)
\item[$\square$]
  No data leakage to validation/test
\item[$\square$]
  Test predictions generated for all samples
\end{itemize}

\begin{center}\rule{0.5\linewidth}{0.5pt}\end{center}

\subsection{WP6: Model Evaluation \&
Selection}\label{wp6-model-evaluation-selection}

\textbf{Duration}: 16--20 June \textbar{} \textbf{Status}: To Do

\subsubsection{Tasks}\label{tasks-4}

\begin{enumerate}
\def\labelenumi{\arabic{enumi}.}
\tightlist
\item
  Cross-validation for all 5 models (k=5 fold)
\item
  Performance benchmarking (RMSE, MAE, R²)
\item
  Create model comparison table and visualizations
\item
  Select best model based on criteria:

  \begin{itemize}
  \tightlist
  \item
    RMSE \textless{} USD 80,000
  \item
    R² \textgreater{} 0.80
  \end{itemize}
\item
  Extract and visualize feature importance
\end{enumerate}

\subsubsection{Deliverables}\label{deliverables-6}

\begin{itemize}
\tightlist
\item
  Model comparison table (CSV)
\item
  Benchmarking plots (PNG/PDF)
\item
  Best model selection report
\item
  Feature importance ranking
\item
  Final model object saved
\end{itemize}

\subsubsection{Quality Criteria}\label{quality-criteria-5}

\begin{itemize}
\tightlist
\item[$\square$]
  All models evaluated on same test set
\item[$\square$]
  Best model meets RMSE \textless{} 80k \& R² \textgreater{} 0.80
  targets
\item[$\square$]
  Feature importance interpreted
\item[$\square$]
  Selection criteria documented
\end{itemize}

\begin{center}\rule{0.5\linewidth}{0.5pt}\end{center}

\subsection{WP7: Deployment \&
Documentation}\label{wp7-deployment-documentation}

\textbf{Duration}: 20--24 June \textbar{} \textbf{Status}: To Do

\subsubsection{Tasks}\label{tasks-5}

\begin{enumerate}
\def\labelenumi{\arabic{enumi}.}
\tightlist
\item
  Build Shiny App (UI + server logic)
\item
  Create prediction endpoint (returns estimate + confidence interval)
\item
  Add feature importance visualization
\item
  Final validation testing
\item
  Prepare presentation slides
\item
  Assemble final report
\item
  Deploy to shinyapps.io or local server
\end{enumerate}

\subsubsection{Deliverables}\label{deliverables-7}

\begin{itemize}
\tightlist
\item
  Deployed Shiny App
\item
  App source code (app.R)
\item
  Presentation slides (PDF)
\item
  Final report (PDF)
\item
  User guide \& model documentation
\item
  All code on GitHub
\end{itemize}

\subsubsection{Quality Criteria}\label{quality-criteria-6}

\begin{itemize}
\tightlist
\item[$\square$]
  Shiny App loads without errors
\item[$\square$]
  Predictions return in \textless5 seconds
\item[$\square$]
  Confidence intervals displayed
\item[$\square$]
  Presentation covers all key findings
\item[$\square$]
  Report submitted via Moodle by deadline
\end{itemize}

\begin{center}\rule{0.5\linewidth}{0.5pt}\end{center}

\section{Dependencies \& Critical
Path}\label{dependencies-critical-path}

{\def\LTcaptype{none} % do not increment counter
\begin{longtable}[]{@{}lll@{}}
\toprule\noalign{}
WP & Depends On & Notes \\
\midrule\noalign{}
\endhead
\bottomrule\noalign{}
\endlastfoot
WP1 & WP0 & Requires Git repo \& Supabase setup \\
WP2 & WP1 & Requires integrated dataset \\
WP3 & WP2 & Requires clean dataset \\
WP4 & WP3 & Requires feature set \\
WP5 & WP4 & Requires preprocessed data \\
WP6 & WP5 & Requires all trained models \\
WP7 & WP6 & Requires selected best model \\
\end{longtable}
}

\textbf{Critical Path Tasks}: 1. WP1: Data integration (26.05) 2. WP2:
Data cleaning (02.06) 3. WP3: Feature engineering (09.06) 4. WP4:
Train/test split (09.06) 5. WP5: Model training (15.06) 6. WP6: Model
selection (20.06) 7. WP7: Shiny App deployment (23.06)

\begin{center}\rule{0.5\linewidth}{0.5pt}\end{center}

\section{Key Milestones}\label{key-milestones}

{\def\LTcaptype{none} % do not increment counter
\begin{longtable}[]{@{}
  >{\raggedright\arraybackslash}p{(\linewidth - 4\tabcolsep) * \real{0.3333}}
  >{\raggedright\arraybackslash}p{(\linewidth - 4\tabcolsep) * \real{0.3333}}
  >{\raggedright\arraybackslash}p{(\linewidth - 4\tabcolsep) * \real{0.3333}}@{}}
\toprule\noalign{}
\begin{minipage}[b]{\linewidth}\raggedright
Date
\end{minipage} & \begin{minipage}[b]{\linewidth}\raggedright
Event
\end{minipage} & \begin{minipage}[b]{\linewidth}\raggedright
Deliverables
\end{minipage} \\
\midrule\noalign{}
\endhead
\bottomrule\noalign{}
\endlastfoot
\textbf{21.05.2026} & Proposal Presentation & Objectives, data
dictionary, organization, development plan \\
\textbf{06.06.2026} & Interim Results & Data status, EDA findings,
initial model results \\
\textbf{18.06.2026} & Feedback Round & Review feedback, adjust if
needed \\
\textbf{24.06.2026} & Final Presentation & Live demo + report submission
via Moodle \\
\end{longtable}
}

\begin{center}\rule{0.5\linewidth}{0.5pt}\end{center}

\section{Weekly Status Checklist}\label{weekly-status-checklist}

\begin{itemize}
\tightlist
\item[$\square$]
  \textbf{Week 1 (13--19.05)}: WP0 setup complete
\item[$\square$]
  \textbf{Week 2 (20--26.05)}: WP1 data integration done, Proposal
  prepared
\item[$\square$]
  \textbf{Week 3 (27.05--02.06)}: WP2 data cleaning complete, WP3
  started
\item[$\square$]
  \textbf{Week 4 (03--09.06)}: WP3 features complete, WP4 preprocessing
  ready
\item[$\square$]
  \textbf{Week 5 (10--16.06)}: WP5 all 5 models trained and validated
\item[$\square$]
  \textbf{Week 6 (17--23.06)}: WP6 best model selected, WP7 Shiny App
  deployed
\item[$\square$]
  \textbf{Week 7 (24.06)}: Final presentation \& report submitted
\end{itemize}

\bookmarksetup{startatroot}

\chapter{Variable selection and Data
Cleaning}\label{variable-selection-and-data-cleaning}

After all the data has been mapped put in a database we construct the
final selection of variables used for modelling and clean this data.

\begin{Shaded}
\begin{Highlighting}[]
\FunctionTok{library}\NormalTok{(DBI)}
\FunctionTok{library}\NormalTok{(RPostgres)}
\end{Highlighting}
\end{Shaded}

\begin{verbatim}
Warning: package 'RPostgres' was built under R version 4.5.3
\end{verbatim}

\begin{Shaded}
\begin{Highlighting}[]
\FunctionTok{library}\NormalTok{(dotenv)}
\end{Highlighting}
\end{Shaded}

\begin{verbatim}
Warning: package 'dotenv' was built under R version 4.5.2
\end{verbatim}

\begin{Shaded}
\begin{Highlighting}[]
\FunctionTok{library}\NormalTok{(here)}
\end{Highlighting}
\end{Shaded}

\begin{verbatim}
here() starts at D:/Documents/00-Personal/01-FHTW/SS2026/MAD-BB-2-SS2026-SOE-R/99-Project-SolutionEngineering/SOE-R-House-Price-Prediction
\end{verbatim}

\begin{Shaded}
\begin{Highlighting}[]
\FunctionTok{library}\NormalTok{(dplyr)}
\end{Highlighting}
\end{Shaded}

\begin{verbatim}

Attaching package: 'dplyr'
\end{verbatim}

\begin{verbatim}
The following objects are masked from 'package:stats':

    filter, lag
\end{verbatim}

\begin{verbatim}
The following objects are masked from 'package:base':

    intersect, setdiff, setequal, union
\end{verbatim}

\begin{Shaded}
\begin{Highlighting}[]
\FunctionTok{library}\NormalTok{(tidyr)}
\FunctionTok{library}\NormalTok{(forcats)}
\FunctionTok{library}\NormalTok{(stringr)}

\FunctionTok{load\_dot\_env}\NormalTok{(here}\SpecialCharTok{::}\FunctionTok{here}\NormalTok{(}\StringTok{".env"}\NormalTok{))}
\end{Highlighting}
\end{Shaded}

\begin{Shaded}
\begin{Highlighting}[]
\NormalTok{con }\OtherTok{\textless{}{-}} \FunctionTok{dbConnect}\NormalTok{(}
\NormalTok{  RPostgres}\SpecialCharTok{::}\FunctionTok{Postgres}\NormalTok{(),}
  \AttributeTok{host     =} \FunctionTok{Sys.getenv}\NormalTok{(}\StringTok{"DB\_HOST"}\NormalTok{),}
  \AttributeTok{port     =} \FunctionTok{as.integer}\NormalTok{(}\FunctionTok{Sys.getenv}\NormalTok{(}\StringTok{"DB\_PORT"}\NormalTok{)),}
  \AttributeTok{dbname   =} \FunctionTok{Sys.getenv}\NormalTok{(}\StringTok{"DB\_NAME"}\NormalTok{),}
  \AttributeTok{user     =} \FunctionTok{Sys.getenv}\NormalTok{(}\StringTok{"DB\_USER"}\NormalTok{),}
  \AttributeTok{password =} \FunctionTok{Sys.getenv}\NormalTok{(}\StringTok{"DB\_PASSWORD"}\NormalTok{)}
\NormalTok{)}
\FunctionTok{dbIsValid}\NormalTok{(con)}
\end{Highlighting}
\end{Shaded}

\begin{verbatim}
[1] TRUE
\end{verbatim}

\section{Subset the Data}\label{subset-the-data}

After mapping the data out and looking at the variable description, we
limit the potentially used features. These features should be ones, that
a homeowner likely knows about their house.

Restrict the features to: - land square feet - waterfront type - view
quality - electricity availability - water availability - street type -
coordinates (longitude, latitude)

\begin{Shaded}
\begin{Highlighting}[]
\NormalTok{df\_1 }\OtherTok{\textless{}{-}} \FunctionTok{dbGetQuery}\NormalTok{(con, }\StringTok{\textquotesingle{}}
\StringTok{  SELECT parcel\_no, appraisal\_account\_type, buildings, land\_net\_squarefeet, waterfront\_type, view\_quality,}
\StringTok{  utility\_electric, utility\_water, street\_type, latitude, longitude}
\StringTok{  FROM "APPRAISAL\_ACCOUNT"}
\StringTok{  WHERE appraisal\_account\_type = }\SpecialCharTok{\textbackslash{}\textquotesingle{}}\StringTok{Residential}\SpecialCharTok{\textbackslash{}\textquotesingle{}}
\StringTok{\textquotesingle{}}\NormalTok{)}
\end{Highlighting}
\end{Shaded}

We get property specifics: - condition - number of stories - number of
bedrooms - number of bathrooms - units - year built - year remodeled

\begin{Shaded}
\begin{Highlighting}[]
\NormalTok{df\_2 }\OtherTok{\textless{}{-}} \FunctionTok{dbGetQuery}\NormalTok{(con, }\StringTok{"}
\StringTok{  SELECT i.parcel\_no, i.building\_id, i.property\_type, i.condition,}
\StringTok{         b.stories, b.bedrooms, b.bathrooms, b.units, b.year\_built, b.year\_remodeled}
\StringTok{  FROM }\SpecialCharTok{\textbackslash{}"}\StringTok{IMPROVEMENT}\SpecialCharTok{\textbackslash{}"}\StringTok{ i}
\StringTok{  LEFT JOIN }\SpecialCharTok{\textbackslash{}"}\StringTok{IMPROVEMENT\_BUILT\_AS}\SpecialCharTok{\textbackslash{}"}\StringTok{ b }
\StringTok{    ON i.parcel\_no = b.parcel\_no }
\StringTok{    AND i.building\_id = b.building\_id}
\StringTok{  WHERE i.property\_type = \textquotesingle{}Residential\textquotesingle{}}
\StringTok{"}\NormalTok{)}
\end{Highlighting}
\end{Shaded}

We get the current property value.

\begin{Shaded}
\begin{Highlighting}[]
\NormalTok{df\_3 }\OtherTok{\textless{}{-}} \FunctionTok{dbGetQuery}\NormalTok{(con, }\StringTok{\textquotesingle{}}
\StringTok{  SELECT parcel\_no, total\_market\_value\_current\_year}
\StringTok{  FROM "TAX\_ACCOUNT"}
\StringTok{\textquotesingle{}}\NormalTok{)}
\end{Highlighting}
\end{Shaded}

we join the data together.

\begin{Shaded}
\begin{Highlighting}[]
\NormalTok{df\_all }\OtherTok{\textless{}{-}}\NormalTok{ df\_1 }\SpecialCharTok{\%\textgreater{}\%}
  \FunctionTok{left\_join}\NormalTok{(df\_2, }\AttributeTok{by =} \StringTok{"parcel\_no"}\NormalTok{) }\SpecialCharTok{\%\textgreater{}\%}
  \FunctionTok{left\_join}\NormalTok{(df\_3, }\AttributeTok{by =} \StringTok{"parcel\_no"}\NormalTok{)}
\end{Highlighting}
\end{Shaded}

\subsection{Cleaning}\label{cleaning}

We obtain the dataset with selected features. Now we go ahead to prepare
the data for initial exploratory data analysis.

\begin{Shaded}
\begin{Highlighting}[]
\NormalTok{df\_all }\OtherTok{\textless{}{-}}\NormalTok{ df\_all }\SpecialCharTok{\%\textgreater{}\%} 
  \FunctionTok{select}\NormalTok{( }\SpecialCharTok{{-}}\NormalTok{property\_type)}
\end{Highlighting}
\end{Shaded}

We first select lands with buildings on it, as we are interested in
residential home valuations.

\begin{Shaded}
\begin{Highlighting}[]
\NormalTok{df\_all }\OtherTok{\textless{}{-}}\NormalTok{ df\_all }\SpecialCharTok{\%\textgreater{}\%} 
  \FunctionTok{filter}\NormalTok{(buildings }\SpecialCharTok{!=} \DecValTok{0}\NormalTok{)}
\end{Highlighting}
\end{Shaded}

Then we look at the missing values and treat them one-by-one.

\begin{Shaded}
\begin{Highlighting}[]
\FunctionTok{colSums}\NormalTok{(}\FunctionTok{is.na}\NormalTok{(df\_all))}
\end{Highlighting}
\end{Shaded}

\begin{verbatim}
                      parcel_no          appraisal_account_type 
                              0                               0 
                      buildings             land_net_squarefeet 
                              0                               3 
                waterfront_type                    view_quality 
                         256540                          252509 
               utility_electric                   utility_water 
                              0                               0 
                    street_type                        latitude 
                              0                               5 
                      longitude                     building_id 
                              5                           29816 
                      condition                         stories 
                          29816                           29846 
                       bedrooms                       bathrooms 
                          29846                           29846 
                          units                      year_built 
                          29887                           29846 
                 year_remodeled total_market_value_current_year 
                          42869                               0 
\end{verbatim}

We remove observations that have no bedroom or bathroom recorded as
these would be one of the features easiest for the user to input and we
want to use them. If they are missing it likely hints at a data
recording issue or a building type that is not a classical residential
home (which is out of scope for this project).

\begin{Shaded}
\begin{Highlighting}[]
\NormalTok{df\_all }\OtherTok{\textless{}{-}}\NormalTok{ df\_all }\SpecialCharTok{\%\textgreater{}\%} 
  \FunctionTok{filter}\NormalTok{(}\SpecialCharTok{!}\FunctionTok{is.na}\NormalTok{(bathrooms))}
\end{Highlighting}
\end{Shaded}

We also remove the 5 observations missing coordinates or square footage.

\begin{Shaded}
\begin{Highlighting}[]
\NormalTok{df\_all }\OtherTok{\textless{}{-}}\NormalTok{ df\_all }\SpecialCharTok{\%\textgreater{}\%} 
  \FunctionTok{filter}\NormalTok{(}\SpecialCharTok{!}\FunctionTok{is.na}\NormalTok{(latitude) }\SpecialCharTok{\&} \SpecialCharTok{!}\FunctionTok{is.na}\NormalTok{(longitude) }\SpecialCharTok{\&} \SpecialCharTok{!}\FunctionTok{is.na}\NormalTok{(land\_net\_squarefeet))}
\end{Highlighting}
\end{Shaded}

Waterfront type records River, Lake, Stream/Creek or Salt. We fill in
missing values to represent no waterfront, as such category does not
exist and missing specification likely represents that the property ahs
no direct water connection.

\begin{Shaded}
\begin{Highlighting}[]
\NormalTok{df\_all }\OtherTok{\textless{}{-}}\NormalTok{ df\_all }\SpecialCharTok{\%\textgreater{}\%} 
  \FunctionTok{mutate}\NormalTok{(}\AttributeTok{waterfront\_type =} \FunctionTok{replace\_na}\NormalTok{(waterfront\_type, }\StringTok{"None"}\NormalTok{))}
\FunctionTok{unique}\NormalTok{(df\_all}\SpecialCharTok{$}\NormalTok{waterfront\_type)}
\end{Highlighting}
\end{Shaded}

\begin{verbatim}
[1] "WF Salt"         "None"            "WF Lake"         "WF Stream/Creek"
[5] "WF River"       
\end{verbatim}

View quality currently has 9 categories that is slightly excessive, and
more confusing for the potential user and increases model complexity.

\begin{Shaded}
\begin{Highlighting}[]
\FunctionTok{unique}\NormalTok{(df\_all}\SpecialCharTok{$}\NormalTok{view\_quality)}
\end{Highlighting}
\end{Shaded}

\begin{verbatim}
 [1] NA              "View Lim"      "View Lim -"    "View Lim +"   
 [5] "View Avg +"    "View Avg"      "View Good"     "View Good +"  
 [9] "View V-Good"   "View V-Good +"
\end{verbatim}

We collapse the categories into ``limited'', ``average'', ``good'',
``very good''. We replace NAs with uncategorized.

\begin{Shaded}
\begin{Highlighting}[]
\NormalTok{df\_all }\OtherTok{\textless{}{-}}\NormalTok{ df\_all }\SpecialCharTok{\%\textgreater{}\%}
  \FunctionTok{mutate}\NormalTok{(}
    \AttributeTok{view\_quality =} \FunctionTok{replace\_na}\NormalTok{(view\_quality, }\StringTok{"uncategorized"}\NormalTok{),}

    \AttributeTok{view\_quality =} \FunctionTok{fct\_collapse}\NormalTok{(view\_quality,}
      \StringTok{"limited"}   \OtherTok{=} \FunctionTok{c}\NormalTok{(}\StringTok{"View Lim"}\NormalTok{, }\StringTok{"View Lim {-}"}\NormalTok{, }\StringTok{"View Lim +"}\NormalTok{),}
      \StringTok{"average"}   \OtherTok{=} \FunctionTok{c}\NormalTok{(}\StringTok{"View Avg"}\NormalTok{, }\StringTok{"View Avg +"}\NormalTok{),}
      \StringTok{"good"}      \OtherTok{=} \FunctionTok{c}\NormalTok{(}\StringTok{"View Good"}\NormalTok{, }\StringTok{"View Good +"}\NormalTok{),}
      \StringTok{"very good"} \OtherTok{=} \FunctionTok{c}\NormalTok{(}\StringTok{"View V{-}Good"}\NormalTok{, }\StringTok{"View V{-}Good +"}\NormalTok{)}
\NormalTok{    )}
\NormalTok{  )}
\FunctionTok{unique}\NormalTok{(df\_all}\SpecialCharTok{$}\NormalTok{view\_quality)}
\end{Highlighting}
\end{Shaded}

\begin{verbatim}
[1] uncategorized limited       average       good          very good    
Levels: uncategorized average good limited very good
\end{verbatim}

We order the view types.

\begin{Shaded}
\begin{Highlighting}[]
\NormalTok{df\_all}\SpecialCharTok{$}\NormalTok{view\_quality }\OtherTok{\textless{}{-}} \FunctionTok{factor}\NormalTok{(df\_all}\SpecialCharTok{$}\NormalTok{view\_quality, }
                               \AttributeTok{levels =} \FunctionTok{c}\NormalTok{(}\StringTok{"uncategorized"}\NormalTok{, }\StringTok{"limited"}\NormalTok{, }\StringTok{"average"}\NormalTok{, }\StringTok{"good"}\NormalTok{, }\StringTok{"very good"}\NormalTok{),}
                               \AttributeTok{ordered =} \ConstantTok{TRUE}\NormalTok{)}
\end{Highlighting}
\end{Shaded}

For the year remodeled, we will proceed as no remodeling happened, if
the year is not recorded.

\begin{Shaded}
\begin{Highlighting}[]
\NormalTok{df\_all }\OtherTok{\textless{}{-}}\NormalTok{ df\_all }\SpecialCharTok{\%\textgreater{}\%} 
  \FunctionTok{mutate}\NormalTok{(}\AttributeTok{year\_remodeled =} \FunctionTok{replace\_na}\NormalTok{(year\_remodeled, }\DecValTok{0}\NormalTok{))}
\end{Highlighting}
\end{Shaded}

Lastly, in the units feature we convert the few missing values to 1 unit
after looking into the data manually, and also convert 0 units to 1 as
the empirical meaning of having no separate units and the whole unit
being one (for example family home) is the same.

\begin{Shaded}
\begin{Highlighting}[]
\NormalTok{df\_all }\OtherTok{\textless{}{-}}\NormalTok{ df\_all }\SpecialCharTok{\%\textgreater{}\%}
  \FunctionTok{mutate}\NormalTok{(}\AttributeTok{units =} \FunctionTok{if\_else}\NormalTok{(}\FunctionTok{is.na}\NormalTok{(units) }\SpecialCharTok{|}\NormalTok{ units }\SpecialCharTok{==} \DecValTok{0}\NormalTok{, }\DecValTok{1}\NormalTok{, units))}
\end{Highlighting}
\end{Shaded}

We have now the complete dataset with the preselected features.

\begin{Shaded}
\begin{Highlighting}[]
\FunctionTok{colSums}\NormalTok{(}\FunctionTok{is.na}\NormalTok{(df\_all))}
\end{Highlighting}
\end{Shaded}

\begin{verbatim}
                      parcel_no          appraisal_account_type 
                              0                               0 
                      buildings             land_net_squarefeet 
                              0                               0 
                waterfront_type                    view_quality 
                              0                               0 
               utility_electric                   utility_water 
                              0                               0 
                    street_type                        latitude 
                              0                               0 
                      longitude                     building_id 
                              0                               0 
                      condition                         stories 
                              0                               0 
                       bedrooms                       bathrooms 
                              0                               0 
                          units                      year_built 
                              0                               0 
                 year_remodeled total_market_value_current_year 
                              0                               0 
\end{verbatim}

A few post checks we should still aplly.

\begin{Shaded}
\begin{Highlighting}[]
\FunctionTok{unique}\NormalTok{(df\_all}\SpecialCharTok{$}\NormalTok{utility\_electric)}
\end{Highlighting}
\end{Shaded}

\begin{verbatim}
[1] "POWER INSTALLED"    "POWER AVAILABLE"    "POWER NO - COMMENT"
\end{verbatim}

\begin{Shaded}
\begin{Highlighting}[]
\NormalTok{df\_all }\OtherTok{\textless{}{-}}\NormalTok{ df\_all }\SpecialCharTok{\%\textgreater{}\%}
  \FunctionTok{mutate}\NormalTok{(}\AttributeTok{utility\_electric =} \FunctionTok{str\_replace}\NormalTok{(utility\_electric, }\StringTok{"POWER NO {-} COMMENT"}\NormalTok{, }\StringTok{"POWER NO"}\NormalTok{))}
\FunctionTok{unique}\NormalTok{(df\_all}\SpecialCharTok{$}\NormalTok{utility\_electric)}
\end{Highlighting}
\end{Shaded}

\begin{verbatim}
[1] "POWER INSTALLED" "POWER AVAILABLE" "POWER NO"       
\end{verbatim}

\begin{Shaded}
\begin{Highlighting}[]
\FunctionTok{unique}\NormalTok{(df\_all}\SpecialCharTok{$}\NormalTok{utility\_water)}
\end{Highlighting}
\end{Shaded}

\begin{verbatim}
[1] "WATER INSTALLED" "WATER NO"        "WATER AVAILABLE"
\end{verbatim}

\begin{Shaded}
\begin{Highlighting}[]
\FunctionTok{unique}\NormalTok{(df\_all}\SpecialCharTok{$}\NormalTok{street\_type)}
\end{Highlighting}
\end{Shaded}

\begin{verbatim}
[1] "STREET UNPAVED" "PAVED"          "STREET NO ROAD"
\end{verbatim}

Check coordinate range.

\begin{Shaded}
\begin{Highlighting}[]
\NormalTok{df\_all }\SpecialCharTok{\%\textgreater{}\%}
  \FunctionTok{summarise}\NormalTok{(}
    \AttributeTok{min\_lat =} \FunctionTok{min}\NormalTok{(latitude, }\AttributeTok{na.rm =} \ConstantTok{TRUE}\NormalTok{),}
    \AttributeTok{max\_lat =} \FunctionTok{max}\NormalTok{(latitude, }\AttributeTok{na.rm =} \ConstantTok{TRUE}\NormalTok{),}
    \AttributeTok{min\_lon =} \FunctionTok{min}\NormalTok{(longitude, }\AttributeTok{na.rm =} \ConstantTok{TRUE}\NormalTok{),}
    \AttributeTok{max\_lon =} \FunctionTok{max}\NormalTok{(longitude, }\AttributeTok{na.rm =} \ConstantTok{TRUE}\NormalTok{)}
\NormalTok{  )}
\end{Highlighting}
\end{Shaded}

\begin{verbatim}
   min_lat  max_lat   min_lon   max_lon
1 46.73973 47.40335 -122.8406 -121.4681
\end{verbatim}

\begin{Shaded}
\begin{Highlighting}[]
\FunctionTok{unique}\NormalTok{(df\_all}\SpecialCharTok{$}\NormalTok{condition)}
\end{Highlighting}
\end{Shaded}

\begin{verbatim}
[1] "Average"       "Extra Poor"    "Uninhabitable" "Poor"         
[5] "Fair"          "Very Poor"    
\end{verbatim}

\begin{Shaded}
\begin{Highlighting}[]
\NormalTok{df\_all}\SpecialCharTok{$}\NormalTok{condition }\OtherTok{\textless{}{-}} \FunctionTok{factor}\NormalTok{(df\_all}\SpecialCharTok{$}\NormalTok{condition, }
                                \AttributeTok{levels =} \FunctionTok{c}\NormalTok{(}\StringTok{"Uninhabitable"}\NormalTok{, }
                                           \StringTok{"Extra Poor"}\NormalTok{, }
                                           \StringTok{"Very Poor"}\NormalTok{, }
                                           \StringTok{"Poor"}\NormalTok{, }
                                           \StringTok{"Fair"}\NormalTok{, }
                                           \StringTok{"Average"}\NormalTok{),}
                                \AttributeTok{ordered =} \ConstantTok{TRUE}\NormalTok{)}
\end{Highlighting}
\end{Shaded}

\begin{Shaded}
\begin{Highlighting}[]
\NormalTok{df\_all }\SpecialCharTok{\%\textgreater{}\%}
  \FunctionTok{select}\NormalTok{(stories, bedrooms, bathrooms) }\SpecialCharTok{\%\textgreater{}\%}  
  \FunctionTok{slice\_max}\NormalTok{(}\AttributeTok{order\_by =}\NormalTok{ stories, }\AttributeTok{n =} \DecValTok{5}\NormalTok{)}
\end{Highlighting}
\end{Shaded}

\begin{verbatim}
  stories bedrooms bathrooms
1      15        3      2.00
2      10        3      2.00
3       8        0      0.50
4       4        3      3.75
5       4        3      4.00
\end{verbatim}

The top 3 observation are very unlikely with this combination of
bedrooms and bathrooms, so we remove them.

\begin{Shaded}
\begin{Highlighting}[]
\NormalTok{to\_remove }\OtherTok{\textless{}{-}}\NormalTok{ df\_all }\SpecialCharTok{\%\textgreater{}\%}
  \FunctionTok{slice\_max}\NormalTok{(}\AttributeTok{order\_by =}\NormalTok{ stories, }\AttributeTok{n =} \DecValTok{3}\NormalTok{, }\AttributeTok{with\_ties =} \ConstantTok{FALSE}\NormalTok{)}

\NormalTok{df\_all }\OtherTok{\textless{}{-}}\NormalTok{ df\_all }\SpecialCharTok{\%\textgreater{}\%}
  \FunctionTok{anti\_join}\NormalTok{(to\_remove)}
\end{Highlighting}
\end{Shaded}

\begin{verbatim}
Joining with `by = join_by(parcel_no, appraisal_account_type, buildings,
land_net_squarefeet, waterfront_type, view_quality, utility_electric,
utility_water, street_type, latitude, longitude, building_id, condition,
stories, bedrooms, bathrooms, units, year_built, year_remodeled,
total_market_value_current_year)`
\end{verbatim}

Looking at the data again manually, and in particular the combination of
bathrooms and bedrooms with teh units we also detect some anomalies that
we remove.

\begin{Shaded}
\begin{Highlighting}[]
\NormalTok{df\_all }\OtherTok{\textless{}{-}}\NormalTok{ df\_all }\SpecialCharTok{\%\textgreater{}\%}
  \FunctionTok{filter}\NormalTok{(bathrooms }\SpecialCharTok{\textless{}=} \DecValTok{17} \SpecialCharTok{\&}\NormalTok{ bedrooms }\SpecialCharTok{\textless{}=} \DecValTok{21}\NormalTok{)}
\end{Highlighting}
\end{Shaded}

The year remodeled being 0 means no remodeling has been recorded on the
house.

\begin{Shaded}
\begin{Highlighting}[]
\NormalTok{df\_all }\SpecialCharTok{\%\textgreater{}\%}
  \FunctionTok{summarise}\NormalTok{(}
    \AttributeTok{min\_built     =} \FunctionTok{min}\NormalTok{(year\_built, }\AttributeTok{na.rm =} \ConstantTok{TRUE}\NormalTok{),}
    \AttributeTok{max\_built     =} \FunctionTok{max}\NormalTok{(year\_built, }\AttributeTok{na.rm =} \ConstantTok{TRUE}\NormalTok{),}
    \AttributeTok{min\_remodeled =} \FunctionTok{min}\NormalTok{(year\_remodeled, }\AttributeTok{na.rm =} \ConstantTok{TRUE}\NormalTok{),}
    \AttributeTok{max\_remodeled =} \FunctionTok{max}\NormalTok{(year\_remodeled, }\AttributeTok{na.rm =} \ConstantTok{TRUE}\NormalTok{)}
\NormalTok{  )}
\end{Highlighting}
\end{Shaded}

\begin{verbatim}
  min_built max_built min_remodeled max_remodeled
1      1861      2026             0          2026
\end{verbatim}

We additionally check if the year remodeled is not 0 it is always more
recent than the year built.

\begin{Shaded}
\begin{Highlighting}[]
\NormalTok{impossible\_dates }\OtherTok{\textless{}{-}}\NormalTok{ df\_all }\SpecialCharTok{\%\textgreater{}\%}
  \FunctionTok{filter}\NormalTok{(year\_remodeled }\SpecialCharTok{\textgreater{}} \DecValTok{0} \SpecialCharTok{\&}\NormalTok{ year\_remodeled }\SpecialCharTok{\textless{}}\NormalTok{ year\_built)}

\FunctionTok{nrow}\NormalTok{(impossible\_dates)}
\end{Highlighting}
\end{Shaded}

\begin{verbatim}
[1] 0
\end{verbatim}

Check the valuation range.

\begin{Shaded}
\begin{Highlighting}[]
\NormalTok{df\_all }\SpecialCharTok{\%\textgreater{}\%}
  \FunctionTok{summarise}\NormalTok{(}
    \AttributeTok{min\_value     =} \FunctionTok{min}\NormalTok{(total\_market\_value\_current\_year, }\AttributeTok{na.rm =} \ConstantTok{TRUE}\NormalTok{),}
    \AttributeTok{max\_value     =} \FunctionTok{max}\NormalTok{(total\_market\_value\_current\_year, }\AttributeTok{na.rm =} \ConstantTok{TRUE}\NormalTok{)}
\NormalTok{  )}
\end{Highlighting}
\end{Shaded}

\begin{verbatim}
  min_value max_value
1      8400  10627600
\end{verbatim}

\begin{Shaded}
\begin{Highlighting}[]
\CommentTok{\#save(df\_all, file = here::here("data", "df\_all.RData"))}
\end{Highlighting}
\end{Shaded}

\bookmarksetup{startatroot}

\chapter*{References}\label{references}
\addcontentsline{toc}{chapter}{References}

\markboth{References}{References}

\protect\phantomsection\label{refs}




\end{document}
